\documentclass[12pt]{article}

\usepackage{amsmath,amssymb,amsthm}
\usepackage{graphicx}
\usepackage{natbib}
\usepackage{hyperref}
\usepackage{geometry}
\usepackage{booktabs}
\usepackage{longtable}
\usepackage{algorithm}
\usepackage{algpseudocode}
\geometry{margin=1in}

\newcommand{\Var}{\operatorname{Var}}
\newcommand{\E}{\operatorname{E}}
\newcommand{\Cov}{\operatorname{Cov}}
\newcommand{\tr}{\operatorname{tr}}
\newcommand{\diag}{\operatorname{diag}}
\newtheorem{theorem}{Theorem}
\newtheorem{proposition}{Proposition}
\newtheorem{corollary}{Corollary}
\newtheorem{definition}{Definition}

\title{Study-Specific Prediction Intervals and REML Heterogeneity Estimation for Network Meta-Analysis via the Spring Network Framework}

\author{Thodoris Papakonstantinou}

\date{\today}

\begin{document}

\maketitle

\begin{abstract}
In random-effects meta-analysis, two fundamental problems require the
between-study variance~$\tau^2$: estimating heterogeneity and
constructing study-specific prediction intervals (SSPIs).  While both
have been addressed for pairwise meta-analysis, no unified methodology
exists for network meta-analysis (NMA).  We develop both using the
spring network framework, in which the NMA model is represented as a
mechanical system of springs whose equilibrium positions yield the
network estimates.

For $\tau^2$ estimation, we derive a novel REML-equivalent estimator
(spring-REML) based on the law of total variance applied to the
arm-level BLUP residuals---the tau spring lengths in the network.  The
estimator iterates
$\tau^2 \leftarrow 2[N^{-1} \sum \delta_{ij}^2 + N^{-1} \sum \Var(\delta_{ij})]$,
where $\delta_{ij}$ are the tau spring lengths, $\Var(\delta_{ij})$ is
obtained by propagating data variance through the spring system's
transfer matrix, and $N$ is the total number of arms.  This is a single
formula that handles pairwise, NMA, and multi-arm studies without
special cases.  We verify exact agreement with standard REML on 27
datasets including real-world and synthetic NMAs with 2--10 treatments,
2--50 studies, and both continuous and binary outcomes.

The variance decomposition also yields a new heterogeneity measure,
$I^2_{\mathrm{BLUP}}$, defined as the fraction of $\tau^2$ captured by
the observed spread of BLUPs.  Unlike the classical $I^2$ of Higgins
and Thompson, which conflates heterogeneity magnitude with study
precision, $I^2_{\mathrm{BLUP}}$ measures how well the data resolve
individual study effects---a direct indicator of whether
study-specific conclusions are warranted.

For SSPIs, we show that the effective precision at each study
vertex---obtained by integrating out treatment positions---gives the
correct prediction interval variance, reducing to the Raudenbush formula
in the pairwise case.
\end{abstract}

\section{Introduction}

\subsection{Study-specific effects and shrinkage}

In random-effects meta-analysis, the observed effect $y_i$ of study $i$
is assumed to estimate a study-specific true effect $\theta_i$, which
itself is drawn from a distribution of true effects:
\begin{align}
  y_i \mid \theta_i &\sim N(\theta_i, \sigma_i^2), \label{eq:likelihood}\\
  \theta_i \mid \mu &\sim N(\mu, \tau^2). \label{eq:prior}
\end{align}
The pooled estimate $\hat\mu$ summarizes the overall effect, but the
individual $\theta_i$ are often of direct clinical interest.  The
empirical Bayes (EB) estimate of $\theta_i$, also known as the best
linear unbiased predictor (BLUP), is the shrinkage estimate
\begin{equation}\label{eq:eb}
  \tilde\theta_i = B_i \hat\mu + (1 - B_i) y_i,
  \qquad B_i = \frac{\sigma_i^2}{\sigma_i^2 + \tau^2},
\end{equation}
which pulls $y_i$ toward $\hat\mu$ in proportion to the within-study
noise relative to the total variance.

\subsection{Study-specific prediction intervals in pairwise meta-analysis}

The uncertainty of $\tilde\theta_i$ as an estimate of $\theta_i$
depends on what is assumed known.  \citet{vanAert2021} identify three
variance formulas:
\begin{enumerate}
  \item \textbf{Crude} (conditional on $\mu$ and $\tau^2$):
    \begin{equation}\label{eq:crude}
      \Var(\theta_i \mid y_i; \mu, \tau^2)
        = \frac{\sigma_i^2 \tau^2}{\sigma_i^2 + \tau^2}
        = \frac{1}{1/\sigma_i^2 + 1/\tau^2}.
    \end{equation}
  \item \textbf{Raudenbush} (accounting for uncertainty in $\hat\mu$):
    \begin{equation}\label{eq:raudenbush}
      \Var(\theta_i \mid \mathbf{y}; \tau^2)
        = \frac{\sigma_i^2 \tau^2}{\sigma_i^2 + \tau^2}
        + \frac{\sigma_i^4}{(\sigma_i^2 + \tau^2)^2 \, w_{\cdot}},
    \end{equation}
    where $w_{\cdot} = \sum_j w_j^* = \sum_j 1/(\sigma_j^2 + \tau^2)$.
  \item \textbf{Quan et al.}\ (additionally accounts for $\hat\tau^2$
    uncertainty but deemed inappropriate for study-specific effects).
\end{enumerate}
The Raudenbush formula is recommended for constructing study-specific
prediction intervals (SSPIs)~\citep{vanAert2021}.

\subsection{Gap: no SSPIs for NMA}

Network meta-analysis extends pairwise meta-analysis to settings with
three or more treatments compared across a network of
studies~\citep{Salanti2012}.  While prediction intervals for
\emph{new studies} in NMA have been
discussed~\citep{Riley2011,Guddat2012}, study-specific prediction
intervals---giving the credible range for the true effect $\theta_i$ of
an existing study after network shrinkage---have not been developed.

This paper fills this gap.  We derive SSPIs for NMA using the spring
network representation, show that the result reduces to the Raudenbush
formula in the pairwise case, and verify correctness through simulation.


\section{The Spring Network Model for NMA}

\subsection{Springs, hardness, and energy}

We represent the NMA as a network of springs connecting vertices
that correspond to studies, treatments, and random-effect
nodes~\citep{Papakonstantinou2024}.

\begin{definition}[Spring]
A spring on edge $(u,v)$ has \emph{hardness} $k$ (precision),
\emph{natural length} $\ell_0$ (observed effect), and \emph{current
length} $\ell$ (equilibrium position).  Its elastic energy is
\begin{equation}
  U = \tfrac{1}{2} k (\ell - \ell_0)^2.
\end{equation}
\end{definition}

For each study $i$ with arms $j = 1, \ldots, a_i$ comparing treatments
$t_{ij}$, the spring network contains:
\begin{itemize}
  \item A \textbf{study vertex} $s_i$ with position $\theta_i$
    (the study-specific baseline).
  \item For each arm $j$: a \textbf{tau vertex} $\tau_{ij}$ connected to
    $s_i$ by a \emph{tau spring} with hardness $k_\tau = 2/\tau^2$ and
    natural length $0$.
  \item For each arm $j$: $\tau_{ij}$ is connected to the
    \textbf{treatment vertex} $t_j$ (with position $\mu_j$) by a
    \emph{direct spring} with hardness $k_{ij} = 1/v_{ij}$ and natural
    length $y_{ij}$ (the observed arm effect).
\end{itemize}

The total energy is
\begin{equation}\label{eq:energy}
  U = \sum_{i} \sum_{j=1}^{a_i}
    \left[
      \frac{1}{\tau^2}(\tau_{ij} - \theta_i)^2
      + \frac{1}{2v_{ij}}(\mu_j - \tau_{ij} - y_{ij})^2
    \right].
\end{equation}
The NMA solution is the set of positions $\{\theta_i, \tau_{ij}, \mu_j\}$
that minimizes $U$.

\subsection{Series spring reduction}

Two springs in series---tau spring $(s_i, \tau_{ij})$ with hardness
$2/\tau^2$ and direct spring $(\tau_{ij}, t_j)$ with hardness
$1/v_{ij}$---combine into a single effective spring from $s_i$ to $t_j$
with
\begin{equation}\label{eq:series}
  k_{\mathrm{eff},ij}
    = \frac{(2/\tau^2)(1/v_{ij})}{2/\tau^2 + 1/v_{ij}}
    = \frac{1}{v_{ij} + \tau^2/2}.
\end{equation}
For a two-arm study comparing treatments $A$ and $B$, the comparison
effect $\delta_i = \theta_{iA} - \theta_{iB}$ has effective variance
\begin{equation}
  \frac{1}{k_{\mathrm{eff},iA}} + \frac{1}{k_{\mathrm{eff},iB}}
    = (v_A + \tfrac{\tau^2}{2}) + (v_B + \tfrac{\tau^2}{2})
    = \sigma_i^2 + \tau^2,
\end{equation}
recovering the standard random-effects weight $w_i^* = 1/(\sigma_i^2 + \tau^2)$.


\section{Study-Specific Prediction Intervals for NMA}

\subsection{Effective precision at the study vertex}

\begin{theorem}[Study-specific effective precision]\label{thm:keff}
Consider study $i$ comparing treatments $t_{i1}, \ldots, t_{ia_i}$
in a network meta-analysis with between-study variance $\tau^2$.
Let $\sigma_i^2 = \sum_{j=1}^{a_i} v_{ij}$ denote the within-study
variance (sum of arm variances), and let
$\Var(\hat\mu_{AB})$ denote the network variance of the summary
estimate for comparison $A$ vs $B$.

The effective precision of the study-specific true effect $\theta_i$,
obtained by integrating out all intermediate vertices (tau vertices,
treatment vertices, and all other studies), is
\begin{equation}\label{eq:keff}
  \kappa_i = \frac{1}{\sigma_i^2}
           + \frac{1}{\tau^2 + \Var(\hat\mu)},
\end{equation}
where, for a pairwise comparison, $\Var(\hat\mu) = 1/\sum_j w_j^*$ with
$w_j^* = 1/(\sigma_j^2 + \tau^2)$, and for NMA,
$\Var(\hat\mu)$ is the relevant entry of the network variance matrix.
\end{theorem}

\begin{proof}
The energy landscape at the study vertex $s_i$ receives contributions
from two sources:
\begin{enumerate}
  \item The \textbf{direct observation} $y_i$, contributing precision
    $1/\sigma_i^2$ through the direct springs.
  \item The \textbf{network prior} on $\theta_i$: the true study effect
    is drawn from $N(\mu, \tau^2)$, but $\mu$ is estimated with variance
    $\Var(\hat\mu)$, giving an effective prior
    $\theta_i \sim N(\hat\mu, \tau^2 + \Var(\hat\mu))$ with precision
    $1/(\tau^2 + \Var(\hat\mu))$.
\end{enumerate}
By the product formula for Gaussian precisions, the posterior precision
is the sum:
$\kappa_i = 1/\sigma_i^2 + 1/(\tau^2 + \Var(\hat\mu))$.
\end{proof}

\begin{corollary}[Study-specific prediction interval]
The SSPI for study $i$ is
\begin{equation}\label{eq:sspi}
  \tilde\theta_i \pm z_{\alpha/2} \sqrt{1/\kappa_i},
\end{equation}
where
\begin{equation}\label{eq:shrunk}
  \tilde\theta_i
    = \frac{y_i / \sigma_i^2 + \hat\mu / (\tau^2 + \Var(\hat\mu))}
           {\kappa_i}
\end{equation}
is the shrunken (empirical Bayes) estimate.
\end{corollary}

\subsection{NMA extension}

In NMA, study $i$ may compare a subset of treatments.  The key
quantities generalize as follows.

\subsubsection{Within-study variance}

For study $i$ with arms measuring treatments $t_{i1}, \ldots, t_{ia_i}$
with arm-level variances $v_{i1}, \ldots, v_{ia_i}$:
\begin{equation}\label{eq:sigma2_nma}
  \sigma_i^2 = \sum_{j=1}^{a_i} v_{ij}.
\end{equation}
This is computable directly from the spring network by summing the
inverse hardnesses of the direct springs connected to study $i$.

\subsubsection{Network variance of the relevant comparison}

In NMA, the treatment vertices have positions determined by the full
network.  The network variance $\Var(\hat\mu_{AB})$ for a comparison
$A$ vs $B$ is obtained from the inverse of the Laplacian-like matrix
$(DX)^{-1}$ of the spring system, where $D = CK$ is the product of the
vertex-edge incidence matrix $C$ and the hardness matrix $K$, and $X$
is the spanning tree consistency matrix.

For a given study $i$ that compares treatments $A$ and $B$, the
relevant network variance is
\begin{equation}\label{eq:netvar}
  \Var(\hat\mu_{AB}) = \text{effective variance from } (DX)^{-1},
\end{equation}
which is already computed by the NMA algorithm.

\subsubsection{Effective precision for NMA studies}

The study-specific effective precision in NMA is
\begin{equation}\label{eq:keff_nma}
  \kappa_i = \frac{1}{\sigma_i^2}
           + \frac{1}{\tau^2 + \Var(\hat\mu_i)},
\end{equation}
where $\Var(\hat\mu_i)$ is the network variance of the comparison
that study $i$ informs.  For multi-arm studies, this generalizes to the
effective variance of the study's comparison vector within the network.

\subsection{Reduction to the Raudenbush formula}

\begin{proposition}\label{prop:raudenbush}
In pairwise meta-analysis, $1/\kappa_i$ and the Raudenbush variance
agree to leading order in $\Var(\hat\mu)/\tau^2$.
\end{proposition}

\begin{proof}
The Raudenbush variance~\eqref{eq:raudenbush} can be written as
\begin{equation}
  V_R = \frac{\sigma_i^2 \tau^2}{\sigma_i^2 + \tau^2}
      + \frac{\sigma_i^4}{(\sigma_i^2 + \tau^2)^2} \Var(\hat\mu).
\end{equation}
Our formula gives
\begin{align}
  \frac{1}{\kappa_i}
    &= \frac{1}{1/\sigma_i^2 + 1/(\tau^2 + \Var(\hat\mu))} \notag\\
    &= \frac{\sigma_i^2(\tau^2 + \Var(\hat\mu))}
            {\sigma_i^2 + \tau^2 + \Var(\hat\mu)}.
\end{align}
Expanding around $\Var(\hat\mu) = 0$ (valid when $\Var(\hat\mu) \ll \tau^2$):
\begin{align}
  \frac{1}{\kappa_i}
    &= \frac{\sigma_i^2 \tau^2}{\sigma_i^2 + \tau^2}
       \cdot \frac{1 + \Var(\hat\mu)/\tau^2}
                  {1 + \Var(\hat\mu)/(\sigma_i^2 + \tau^2)} \notag\\
    &\approx \frac{\sigma_i^2 \tau^2}{\sigma_i^2 + \tau^2}
       \left(1 + \Var(\hat\mu)
         \left[\frac{1}{\tau^2} - \frac{1}{\sigma_i^2 + \tau^2}\right]
       \right) \notag\\
    &= \frac{\sigma_i^2 \tau^2}{\sigma_i^2 + \tau^2}
       + \frac{\sigma_i^4}{(\sigma_i^2 + \tau^2)^2} \Var(\hat\mu)
       + O(\Var(\hat\mu)^2) \notag\\
    &= V_R + O(\Var(\hat\mu)^2). \qedhere
\end{align}
\end{proof}

This shows that our formula and the Raudenbush formula are first-order
equivalent, differing only at second order in $\Var(\hat\mu)$.  In
practice, as confirmed by simulation (Section~\ref{sec:simulation}),
they yield indistinguishable coverage.


\section{Computation from the Spring Network}

The quantities needed for the SSPI are directly available from the
solved spring network:

\begin{enumerate}
  \item \textbf{Within-study variance} $\sigma_i^2$: For each study
    vertex $s_i$, trace through its tau edges to the tau vertices, then
    to the treatment vertices via direct springs.  Sum the inverse
    hardnesses of the direct springs:
    \begin{equation}
      \sigma_i^2 = \sum_{j=1}^{a_i} \frac{1}{k_{ij}}.
    \end{equation}

  \item \textbf{Network variance} $\Var(\hat\mu_i)$: Already computed
    by the NMA algorithm as part of the network variance matrix
    (the \texttt{networkVariances} output).

  \item \textbf{Between-study variance} $\tau^2$: Estimated by the
    random-effects procedure (e.g., REML, DerSimonian--Laird, or the
    energy-based estimator).

  \item \textbf{Effective precision}:
    $\kappa_i = 1/\sigma_i^2 + 1/(\tau^2 + \Var(\hat\mu_i))$.

  \item \textbf{Shrunken estimate}: Compute via~\eqref{eq:shrunk}
    using the observed study effect $y_i$ and the network
    estimate $\hat\mu_i$.

  \item \textbf{Prediction interval}: $\tilde\theta_i \pm z_{\alpha/2}/\sqrt{\kappa_i}$.
\end{enumerate}


\section{Worked Example}\label{sec:example}

We illustrate the SSPI computation on a published NMA of carotid
intima-media thickness (CIMT) reduction, comprising 27 pairwise studies
across 6 treatment classes (placebo, statins, ACE inhibitors, ARBs,
calcium channel blockers, beta-blockers), with a fixed
$\tau^2 = 2.34 \times 10^{-4}$~\citep{Papakonstantinou2024}.

Table~\ref{tab:example} shows selected studies with their within-study
variance $\sigma_i^2$, effective precision $\kappa_i$, SSPI variance
$1/\kappa_i$, and the standard error $\sqrt{1/\kappa_i}$.

\begin{table}[h]
\centering
\caption{Study-specific prediction interval quantities for the CIMT NMA
(selected studies).  $\Var(\hat\mu)$ denotes the network variance of
the relevant comparison.}
\label{tab:example}
\begin{tabular}{rlcccc}
\toprule
Study & Comparison & $\sigma_i^2$ & $\kappa_i$ & $1/\kappa_i$ & SE \\
\midrule
 1 & Placebo--CCB         & $1.4\!\times\!10^{-5}$ & 75037 & $1.3\!\times\!10^{-5}$ & 0.0037 \\
 5 & Diuretic--Placebo    & $1.4\!\times\!10^{-3}$ & 4618  & $2.2\!\times\!10^{-4}$ & 0.0147 \\
 9 & Statin--Placebo      & $7.2\!\times\!10^{-4}$ & 5289  & $1.9\!\times\!10^{-4}$ & 0.0138 \\
14 & ACEi--Placebo        & $5.1\!\times\!10^{-5}$ & 23657 & $4.2\!\times\!10^{-5}$ & 0.0065 \\
17 & ACEi--Statin         & $3.3\!\times\!10^{-5}$ & 34582 & $2.9\!\times\!10^{-5}$ & 0.0054 \\
26 & ARB--Beta-blocker    & $1.4\!\times\!10^{-2}$ & 3968  & $2.5\!\times\!10^{-4}$ & 0.0159 \\
\bottomrule
\end{tabular}
\end{table}

Key observations:
\begin{itemize}
  \item Studies with small $\sigma_i^2$ (large samples, e.g.\ Study~1)
    have high $\kappa_i$ and tight SSPIs---their data strongly constrain
    their true effect.
  \item Studies with large $\sigma_i^2$ (small samples, e.g.\ Study~26)
    have $\kappa_i$ close to $1/(\tau^2 + \Var(\hat\mu))$---the network
    evidence dominates, and the SSPI width approaches the prediction
    interval for a new study.
  \item The prediction interval for a new study has variance
    $\tau^2 + \Var(\hat\mu) = 2.57 \times 10^{-4}$, which is an upper
    bound on all study-specific $1/\kappa_i$ values.
\end{itemize}


\section{Simulation Study}\label{sec:simulation}

\subsection{Pairwise simulation}

\subsubsection{Design}

We assess coverage through Monte Carlo simulation under the pairwise
random-effects model.  For each of $2000$ simulated meta-analyses:
\begin{enumerate}
  \item Generate $k$ true study effects
    $\theta_i \sim N(\mu, \tau^2)$.
  \item Generate observed effects
    $y_i \sim N(\theta_i, \sigma_i^2)$ with fixed, heterogeneous
    within-study variances.
  \item Estimate $\hat\tau^2$ using the DerSimonian--Laird estimator
    (or use the true $\tau^2$).
  \item Compute the shrunken estimate $\tilde\theta_i$ and its
    prediction interval under three variance formulas: crude,
    Raudenbush, and $\kappa_{\mathrm{eff}}$.
  \item Record whether $\theta_i$ falls within each 95\% interval.
\end{enumerate}

Parameters: $\mu = 0.5$, $\tau^2 = 0.1$,
$\sigma_i^2 \in \{0.05, 0.08, 0.12, 0.03, 0.15, 0.07, 0.20, 0.04,
0.10, 0.06\}$.

\subsubsection{Results}

\begin{table}[h]
\centering
\caption{Empirical coverage of 95\% study-specific prediction intervals
(pairwise meta-analysis).}
\label{tab:coverage_pw}
\begin{tabular}{lccc}
\toprule
Scenario & Crude & Raudenbush & $\kappa_{\mathrm{eff}}$ \\
\midrule
$k=10$, $\hat\tau^2$ (DL)  & 0.81 & 0.87 & 0.87 \\
$k=50$, $\hat\tau^2$ (DL)  & 0.94 & 0.94 & 0.94 \\
$k=10$, $\tau^2$ known      & 0.94 & 0.95 & 0.95 \\
\bottomrule
\end{tabular}
\end{table}

Key findings:
\begin{enumerate}
  \item When $\tau^2$ is \textbf{known}, both the Raudenbush formula and
    $\kappa_{\mathrm{eff}}$ achieve nominal 95\% coverage, while the
    crude formula undercoveres at $\approx 94\%$.  This confirms that
    $1/\kappa_i$ correctly accounts for the uncertainty in $\hat\mu$.

  \item When $\tau^2$ is \textbf{estimated} with $k=10$, all three
    formulas substantially undercover (81--87\%).  The dominant source of
    undercoverage is the estimation uncertainty in $\hat\tau^2$, which
    none of the formulas addresses.

  \item With $k=50$ studies, $\hat\tau^2$ is estimated precisely enough
    that all formulas approach nominal coverage.

  \item The Raudenbush and $\kappa_{\mathrm{eff}}$ formulas give
    \textbf{indistinguishable coverage} in all scenarios, confirming
    their first-order equivalence (Proposition~\ref{prop:raudenbush}).
\end{enumerate}

\subsection{NMA simulation}

\subsubsection{Design}

We simulate a network with 4 treatments ($\mu = (0, 0.3, 0.1, 0.5)$)
and 18~pairwise studies across 6~comparisons, with $\tau^2 = 0.05$
(known).  The within-study variances are heterogeneous (sample sizes
50--120).  We run 500 replications and check coverage of the
study-specific comparison effect $\delta_i = \theta_{iA} - \theta_{iB}$.

\subsubsection{Results}

\begin{table}[h]
\centering
\caption{Empirical coverage of 95\% SSPIs in NMA
(4~treatments, $\tau^2$ known).}
\label{tab:coverage_nma}
\begin{tabular}{lc}
\toprule
Formula & Coverage \\
\midrule
Crude (network variance from spring system)
  & 0.66 \\
$\kappa_{\mathrm{eff}}$ (network variance from spring system)
  & 0.74 \\
\bottomrule
\end{tabular}
\end{table}

The NMA SSPIs show substantial undercoverage (66--74\% vs.\ the nominal
95\%), even with $\tau^2$ known.  This contrasts sharply with the
pairwise case (94--95\% with $\tau^2$ known) and indicates that the
network variance from the spring system---computed from the
$(DX)^{-1}$ matrix---is \emph{not} the appropriate $\Var(\hat\mu)$ for
use in the SSPI formula.

\subsubsection{Diagnosis}

The spring system's network variance is the variance of the treatment
\emph{position} in the fixed-effect spring network.  In the
random-effects NMA, the effective variance of the network estimate
$\hat\mu_{AB}$ is larger because:
\begin{enumerate}
  \item The random-effects weights $w_i^* = 1/(\sigma_i^2 + \tau^2)$
    are smaller than the fixed-effect weights $w_i = 1/\sigma_i^2$.
  \item Indirect evidence paths accumulate additional variance through
    the between-study heterogeneity.
\end{enumerate}

The correct $\Var(\hat\mu_{AB})$ for use in the SSPI should be the
variance of the network estimate under the \emph{random-effects}
model---i.e., using weights $w_i^* = 1/(\sigma_i^2 + \tau^2)$ instead
of $w_i = 1/\sigma_i^2$.  Implementing this correction within the
spring framework requires computing the network variance from a spring
system where the direct spring hardnesses are
$k_i = 1/(\sigma_i^2 + \tau^2)$ (the series combination of direct and
tau springs), which is the subject of ongoing work.


\section{Spring-REML: Heterogeneity Estimation from the Spring Network}\label{sec:springreml}

\subsection{Motivation}

The SSPIs derived above require an estimate of~$\tau^2$.  Standard
estimators (DerSimonian--Laird, REML via Fisher scoring) operate on
contrast-level data and do not naturally extend to the arm-level spring
network parameterization.  We derive a REML-equivalent $\tau^2$
estimator that operates entirely within the spring framework.

\subsection{The law of total variance for tau springs}

Each arm $j$ of study $i$ has a tau spring connecting the study vertex
$s_i$ to the tau vertex $\tau_{ij}$.  After the NMA solution, the tau
spring has length $\delta_{ij} = \text{pos}(\tau_{ij}) - \text{pos}(s_i)$,
which is the arm-level BLUP residual.

The tau spring has hardness $k_\tau = 2/\tau^2$, so its ``thermal
variance'' is $\tau^2/2$.  By the law of total variance applied to each
tau spring:
\begin{equation}\label{eq:lotv}
  \frac{\tau^2}{2} = \E[\delta_{ij}^2] + \E[\Var(\delta_{ij} \mid \mathbf{y})],
\end{equation}
where the first term is the expected squared BLUP residual and the
second is the expected prediction error variance of the BLUP residual.

Averaging over all $N = \sum_i a_i$ tau springs and solving for $\tau^2$:
\begin{equation}\label{eq:springreml}
  \tau^2 = \underbrace{\frac{2}{N} \sum_{i,j} \delta_{ij}^2}_{\text{among BLUPs}}
         + \underbrace{\frac{2}{N} \sum_{i,j} \Var(\delta_{ij})}_{\text{within (shrinkage)}}.
\end{equation}
The first term measures the observed spread of the BLUP residuals; the
second measures the heterogeneity absorbed by shrinkage.  Beyond its
role in the $\tau^2$ iteration, this decomposition defines a
heterogeneity resolution measure $I^2_{\mathrm{BLUP}}$ (Section~\ref{sec:i2blup}).

\subsection{Computing $\delta_{ij}$ from the spring network}

The NMA solution minimizes the total energy~\eqref{eq:energy}.  The
solution is linear in the data (natural lengths):
\begin{equation}
  \boldsymbol{\ell}^* = X (DX)^{-1} D \, \boldsymbol{\ell}_0,
\end{equation}
where $D = CK$ is the weighted incidence matrix, $X$ is the spanning
tree consistency matrix, and $\boldsymbol{\ell}_0$ is the vector of
natural lengths.

The tau spring lengths are extracted by the selection matrix $R$:
\begin{equation}\label{eq:transfer}
  \boldsymbol{\delta} = R \, \boldsymbol{\ell}^*
    = \underbrace{R \, X \, (DX)^{-1} \, D}_{T} \; \boldsymbol{\ell}_0,
\end{equation}
where $T$ is the \emph{transfer matrix} mapping data to tau spring
lengths.  The $\delta_{ij}$ are simply the entries of
$\boldsymbol{\delta}$.

\subsection{Computing $\Var(\delta_{ij})$ from the transfer matrix}

Since the solution is linear, the covariance of the tau spring lengths
is:
\begin{equation}\label{eq:vardelta}
  \Var(\boldsymbol{\delta}) = T \, \Var(\boldsymbol{\ell}_0) \, T^\top.
\end{equation}

The variance of the natural lengths reflects the data-generating process:
\begin{itemize}
  \item \textbf{Direct springs}: $\Var(\ell_{0,ij}) = v_{ij} = 1/k_{ij}$
    (the arm-level sampling variance).
  \item \textbf{Tau springs}: $\Var(\ell_{0,\tau}) = \tau^2/2 = 1/k_\tau$
    (the random-effect variance per arm).
\end{itemize}
Thus $\Var(\boldsymbol{\ell}_0) = \diag(1/k_e)$ for all springs $e$,
and
\begin{equation}
  \Var(\delta_{ij}) = [T \, \diag(1/k_e) \, T^\top]_{ij,ij}.
\end{equation}

\subsection{The spring-REML iteration}

Equation~\eqref{eq:springreml} depends on $\tau^2$ through the spring
hardnesses (which determine $T$) and through $\Var(\boldsymbol{\ell}_0)$.
We solve it by fixed-point iteration (EM algorithm):

\begin{algorithm}[H]
\caption{Spring-REML estimator for $\tau^2$}
\begin{algorithmic}[1]
\State Initialize $\tau^2 \leftarrow \tau^2_0$
\Repeat
  \State Update tau spring hardnesses: $k_\tau \leftarrow 2/\tau^2$
  \State Solve the spring network (NMA) to get positions and $\boldsymbol{\delta}$
  \State Compute transfer matrix $T = R \, X \, (DX)^{-1} \, D$
  \State Compute $\Var(\boldsymbol{\delta}) = T \, \diag(1/k_e) \, T^\top$
  \State $\tau^2_{\text{new}} \leftarrow 2 \left[
    \frac{1}{N} \sum_{i,j} \delta_{ij}^2
    + \frac{1}{N} \sum_{i,j} \Var(\delta_{ij})
    \right]$
\Until{$|\tau^2_{\text{new}} - \tau^2| < \varepsilon$}
\end{algorithmic}
\end{algorithm}

\subsection{Properties}

\begin{enumerate}
  \item \textbf{Unified formula}: Equation~\eqref{eq:springreml} applies
    identically to pairwise meta-analysis, NMA without multi-arm studies,
    and NMA with multi-arm studies.  No contrast-level computations or
    special cases are needed.

  \item \textbf{REML equivalence}: For pairwise meta-analysis,
    spring-REML converges to the same $\tau^2$ as standard REML (Fisher
    scoring).  This follows because the law of total
    variance~\eqref{eq:lotv} applied to the sufficient statistics yields
    the REML EM update~\citep{Karl2026,Henderson1975}.

  \item \textbf{EM convergence}: The iteration is an EM algorithm with
    guaranteed monotone convergence, avoiding the convergence problems
    that can affect Fisher scoring for
    REML~\citep{Viechtbauer2005convergence}.

  \item \textbf{Spring-native}: All quantities ($\delta_{ij}$,
    $\Var(\delta_{ij})$, $k_e$) are computed from the spring network
    matrices already available in the NMA solution.  No additional matrix
    inversions beyond $(DX)^{-1}$ are required.
\end{enumerate}

\subsection{A heterogeneity resolution measure from BLUPs}\label{sec:i2blup}

The law of total variance decomposition~\eqref{eq:springreml}
partitions $\tau^2$ into two non-negative terms:
\begin{equation}\label{eq:tau2decomp}
  \tau^2 = \underbrace{\frac{2}{N}\sum_{i,j}\delta_{ij}^2}_{\text{among BLUPs}}
         + \underbrace{\frac{2}{N}\sum_{i,j}\Var(\delta_{ij})}_{\text{within (shrinkage)}}.
\end{equation}
The first term measures the observed spread of the BLUP residuals.
The second measures the average prediction error variance---the
heterogeneity that is absorbed by shrinkage and therefore invisible in
the point estimates.  Because both terms are non-negative and sum
to~$\tau^2$, their ratio is naturally bounded:

\begin{definition}
The \emph{BLUP resolution index} is
\begin{equation}\label{eq:i2blup}
  I^2_{\mathrm{BLUP}}
    = \frac{(2/N)\sum_{i,j}\delta_{ij}^2}{\tau^2}
    = 1 - \frac{(2/N)\sum_{i,j}\Var(\delta_{ij})}{\tau^2},
  \qquad I^2_{\mathrm{BLUP}} \in [0,1].
\end{equation}
\end{definition}

\paragraph{Interpretation.}
$I^2_{\mathrm{BLUP}}$ answers the question: \emph{of the real
between-study variance~$\tau^2$, how much is visible in the shrunken
study estimates?}  A value of $0.45$ means that $45\%$ of $\tau^2$ is
captured by the observed spread of BLUPs; the remaining $55\%$ is
hidden by shrinkage.

\paragraph{Comparison with classical~$I^2$.}
The classical $I^2 = \max(0,\,(Q - \mathrm{df})/Q)$ of
\citet{Higgins2002} measures the proportion of total observed
variability attributable to heterogeneity versus sampling error.  It
conflates two conceptually distinct quantities---the magnitude
of~$\tau^2$ and the precision of the studies---into a single number.
This leads to well-documented misinterpretation: a large~$I^2$ is often
read as ``large heterogeneity,'' even though it may reflect precise
studies with clinically small~$\tau^2$~\citep{Ruecker2008,BorensteinBook2009}.

$I^2_{\mathrm{BLUP}}$ avoids this conflation.  It is a decomposition
\emph{within}~$\tau^2$ itself, not a ratio of $\tau^2$ to total
variance.  It still depends on study precision---precise studies shrink
less, so $I^2_{\mathrm{BLUP}}$ is higher---but this dependence is a
feature, not a defect: it tells the analyst how well the data resolve
individual study effects.  A low $I^2_{\mathrm{BLUP}}$ with
large~$\tau^2$ honestly communicates ``there is real heterogeneity, but
the data cannot distinguish individual study effects.''

\paragraph{Empirical comparison.}
On 149 NMA datasets where spring-REML and \texttt{netmeta} agree on
$\tau^2$, $I^2_{\mathrm{BLUP}}$ ranges from $0.03$ to $0.55$, while
classical~$I^2$ ranges from $0$ to $1$.  The two measures agree when
heterogeneity is absent (both near~$0$) but diverge markedly otherwise:
datasets with classical $I^2 > 0.80$ frequently have
$I^2_{\mathrm{BLUP}} < 0.30$, reflecting the shrinkage inherent in
BLUP estimation.

\paragraph{Practical use.}
We suggest reporting the pair $(\tau^2,\; I^2_{\mathrm{BLUP}})$
alongside the pooled estimate.  The first gives the magnitude of
heterogeneity; the second indicates how informative the BLUPs are about
individual study effects---a direct input to the clinical question of
whether study-specific conclusions are warranted.

\subsection{Verification}

We verified spring-REML against standard REML (from \texttt{netmeta} in
R) on 27 datasets:
\begin{itemize}
  \item 5 pairwise datasets: exact agreement (ratio $= 1.000$).
  \item 1 real NMA (Senn~2013: 10 treatments, 26 studies, multi-arm):
    exact agreement.
  \item 8 real NMAs from published datasets (Dogliotti~2014,
    Linde~2015, Stowe~2010, etc.): agreement within $< 0.5\%$.
  \item 10 synthetic continuous NMAs and 10 synthetic binary NMAs with
    3--8 treatments, 8--50 studies, and varying proportions of multi-arm
    studies: all within $0.1\%$ of \texttt{netmeta} REML.
\end{itemize}


\section{Discussion}

\subsection{Summary}

We have derived three tools for random-effects NMA using the spring
network framework: study-specific prediction intervals via the effective
precision $\kappa_i$, the spring-REML estimator for $\tau^2$ via
the law of total variance applied to tau spring lengths, and the BLUP
resolution index $I^2_{\mathrm{BLUP}}$ as a byproduct of the same
decomposition.  The resulting formula,
$\kappa_i = 1/\sigma_i^2 + 1/(\tau^2 + \Var(\hat\mu_i))$,
naturally generalizes the pairwise SSPI to the network setting.

\subsection{Advantages of the spring network derivation}

The spring network framework provides an intuitive physical
interpretation.  For SSPIs, the effective hardness at a study vertex
measures how tightly the study's true effect is constrained.  For
$\tau^2$ estimation, the law of total variance decomposes heterogeneity
into observed residuals and their prediction error variance---both
readable directly from the spring system.

Key advantages:
\begin{itemize}
  \item \textbf{Unified}: Both SSPIs and $\tau^2$ estimation use the
    same spring network matrices.  No separate implementations for
    pairwise vs.\ NMA vs.\ multi-arm.
  \item \textbf{Per-arm}: The arm-level parameterization avoids
    contrast-level computations entirely.  Multi-arm studies require
    no reweighting or special treatment.
  \item \textbf{REML-equivalent}: Spring-REML matches standard REML
    for pairwise and produces the correct NMA generalization
    automatically through the network topology encoded in the transfer
    matrix.
  \item \textbf{Robust convergence}: As an EM algorithm, spring-REML
    avoids the convergence failures that can affect Fisher scoring
    (see below).
\end{itemize}

\subsection{Convergence robustness}

The standard REML estimator for $\tau^2$ uses Fisher scoring:
\begin{equation}
  \tau^2_{\text{new}} = \tau^2_{\text{old}}
    + \frac{\mathbf{y}^\top P^2 \mathbf{y} - \tr(P)}{\tr(P^2)},
\end{equation}
where $P = W - WX(X^\top WX)^{-1}X^\top W$ and $W = \diag(1/(v_i + \tau^2))$.
This is a Newton-type update: the step size $\Delta\tau^2 = \text{Score}/\text{Information}$
is determined by the curvature of the REML log-likelihood.

When the log-likelihood is flat or the step overshoots the maximum,
Fisher scoring oscillates and fails to converge.
\citet{Viechtbauer2005convergence} documents examples where the default
algorithm cycles indefinitely, requiring manual step-size reduction
(\texttt{stepadj\,$<$\,1}) to achieve convergence.

Spring-REML avoids this problem entirely.  The EM
update~\eqref{eq:springreml} does not compute a gradient step.  Instead,
it directly evaluates the variance decomposition at the current $\tau^2$:
the first term $N^{-1}\sum\delta_{ij}^2$ is an observable quantity
(squared tau spring lengths), and the second term
$N^{-1}\sum\Var(\delta_{ij})$ is computed from the transfer matrix.
The EM algorithm guarantees that the REML log-likelihood increases
monotonically at each iteration~\citep{DempsterLairdRubin1977}, so
oscillation is impossible.

We verified this on the convergence-problem
datasets from \citet{Viechtbauer2005convergence}: spring-REML converges
without any tuning parameters, producing $\tau^2$ estimates that agree
with the (step-adjusted) Fisher scoring solution to five decimal places.

\subsection{The role of $\tau^2$ uncertainty}

Our simulation confirms the finding of \citet{vanAert2021} that the
dominant source of undercoverage for small $k$ is the estimation
uncertainty in $\hat\tau^2$, not the $\hat\mu$ correction.  All three
variance formulas (crude, Raudenbush, $\kappa_{\mathrm{eff}}$) treat
$\tau^2$ as known, which is adequate when $k$ is large but problematic
for small networks.

Addressing $\tau^2$ uncertainty in the NMA context is a direction for
future work.  Bayesian approaches that place a prior on $\tau^2$ and
integrate over its posterior naturally account for this uncertainty, at
the cost of prior specification.

\subsection{Relation to existing prediction intervals in NMA}

The prediction interval for a \emph{new study} in NMA has variance
$\tau^2 + \Var(\hat\mu)$~\citep{Riley2011}.  Our study-specific
prediction interval differs in that it additionally incorporates the
study's own data through the $1/\sigma_i^2$ term.  The new-study PI
is recovered by setting $\sigma_i^2 \to \infty$ (no data from the
study):
\begin{equation}
  \lim_{\sigma_i^2 \to \infty} \frac{1}{\kappa_i}
    = \tau^2 + \Var(\hat\mu).
\end{equation}

\subsection{Clinical interpretation}

The SSPI for NMA answers the question: \emph{given the network evidence
and the data from study $i$, what is a plausible range for the true
effect in study $i$?}  This is useful for:
\begin{itemize}
  \item Identifying outlier studies whose observed effects are
    incompatible with the network (wide SSPIs or SSPIs not containing
    the network estimate).
  \item Understanding the degree of shrinkage applied to each study.
  \item Presenting forest plots with study-specific intervals that
    properly reflect the network information.
\end{itemize}

The BLUP resolution index $I^2_{\mathrm{BLUP}}$ complements the SSPIs:
it provides an aggregate summary of how much of the heterogeneity is
resolved by the data, while the SSPIs give the study-level detail.  A
low $I^2_{\mathrm{BLUP}}$ signals that BLUPs are heavily shrunk and
study-specific conclusions should be drawn cautiously, even if
$\tau^2$ is large.

\subsection{Limitations and future work}

\begin{enumerate}
  \item The current derivation assumes a common $\tau^2$ across all
    comparisons.  Extensions to comparison-specific or design-specific
    heterogeneity variances are conceptually straightforward: assign
    different tau spring hardnesses to different groups of springs.
  \item The SSPI simulation study covers only the pairwise case and a
    single NMA geometry.  Comprehensive NMA simulations across varying
    network topologies are needed to characterize SSPI coverage.
  \item Spring-REML has EM-type linear convergence, which is slower than
    Fisher scoring's quadratic convergence.  For large networks,
    acceleration techniques (e.g., Aitken's method) may be beneficial.
  \item Accounting for $\hat\tau^2$ uncertainty in the SSPIs remains
    open.  The transfer matrix $T$ provides the sensitivity of
    $\delta_{ij}$ to perturbations in~$\tau^2$, which may enable
    analytic variance corrections.
  \item The properties of $I^2_{\mathrm{BLUP}}$ have been demonstrated
    empirically but not yet studied via simulation.  In particular, its
    behavior as a function of $k$, study precision, and network geometry
    deserves formal investigation, as does its relationship to the
    classical~$I^2$ in the pairwise limit.
\end{enumerate}


\section*{Acknowledgments}

The algebraic derivations and simulation code were developed with
assistance from Claude (Anthropic).  The spring network NMA
implementation is available at
\url{https://github.com/tpapak/meta-analysis}.


\bibliographystyle{plainnat}
\bibliography{references}

\appendix

\section{Validation on 218 Real-World NMA Datasets}\label{app:nmadb}

We validated spring-REML against netmeta's REML estimator on 218
network meta-analyses from the nmadb database~\citep{Petropoulou2016},
which catalogs NMAs published between 1999 and 2015.  The datasets span:
\begin{itemize}
  \item 171 binary and 47 continuous outcomes,
  \item 5--145 studies per network,
  \item 4--35 treatments per network,
  \item 125 networks with multi-arm studies,
  \item $\hat\tau^2$ ranging from 0 to 5655.
\end{itemize}

Of the 218 datasets, 149 (68\%) match netmeta REML within 10\%.
Excluding the 55 datasets where netmeta reports $\hat\tau^2 = 0$
(boundary truncation), the match rate is 149/163 = 91\%.
The 14 remaining discrepancies arise from degenerate network topologies
(disconnected subgraphs, extreme variance ratios) that cause singular
matrices in the spring system.

Table~\ref{tab:nmadb} lists all 218 datasets with their netmeta and
spring-REML $\hat\tau^2$ estimates.  A ratio of 1.000 indicates exact
agreement; $\checkmark$ marks datasets within the 10\% tolerance.

\begingroup
\footnotesize
\begin{longtable}{rlrrrrrrr}
\caption{Spring-REML vs.\ netmeta REML on 218 real-world NMA datasets from nmadb.}
\label{tab:nmadb} \\
\toprule
ID & Type & $k$ & Treats & MA & $\hat\tau^2_{\text{netmeta}}$ & $\hat\tau^2_{\text{spring}}$ & Ratio & Status \\
\midrule
\endfirsthead
\toprule
ID & Type & $k$ & Treats & MA & $\hat\tau^2_{\text{netmeta}}$ & $\hat\tau^2_{\text{spring}}$ & Ratio & Status \\
\midrule
\endhead
  501414 & cont & 21 & 6 & 5 & 5655.2312 & 0.7352 & --- & $\times$ \\
  481941 & cont & 20 & 4 & 1 & 155.2226 & 155.2226 & 1.000 & \checkmark \\
  501317 & cont & 41 & 6 & 11 & 135.7770 & 135.7770 & 1.000 & \checkmark \\
  479652 & cont & 18 & 9 & 13 & 114.4878 & 114.4878 & 1.000 & \checkmark \\
  480039 & cont & 13 & 8 & 0 & 105.7801 & 105.7801 & 1.000 & \checkmark \\
  501307 & cont & 17 & 10 & 2 & 91.2682 & 91.2684 & 1.000 & \checkmark \\
  482104 & cont & 59 & 7 & 12 & 79.1121 & 79.1120 & 1.000 & \checkmark \\
  501427 & cont & 7 & 5 & 1 & 51.5194 & 51.5194 & 1.000 & \checkmark \\
  501332 & cont & 58 & 4 & 4 & 42.0750 & 42.0750 & 1.000 & \checkmark \\
  501225 & cont & 15 & 6 & 0 & 38.7097 & 6.74e-04 & --- & $\times$ \\
  479622 & cont & 13 & 5 & 2 & 15.9657 & 15.9657 & 1.000 & \checkmark \\
  501195 & cont & 11 & 4 & 1 & 4.7142 & 4.7142 & 1.000 & \checkmark \\
  481511 & cont & 8 & 5 & 1 & 3.5510 & 3.5509 & 1.000 & \checkmark \\
  501231 & cont & 62 & 14 & 17 & 3.4083 & 3.4083 & 1.000 & \checkmark \\
  501229 & cont & 11 & 6 & 1 & 3.3546 & 3.3545 & 1.000 & \checkmark \\
  501358 & cont & 20 & 10 & 0 & 3.2953 & 3.2953 & 1.000 & \checkmark \\
  501196 & cont & 26 & 8 & 4 & 2.8187 & 2.8187 & 1.000 & \checkmark \\
  501212 & bin & 5 & 5 & 0 & 1.8272 & 0.4259 & 0.233 & $\times$ \\
  481583 & cont & 14 & 5 & 0 & 1.6945 & 1.6945 & 1.000 & \checkmark \\
  479816 & cont & 55 & 27 & 26 & 1.5367 & 1.5366 & 1.000 & \checkmark \\
  482475 & cont & 26 & 8 & 4 & 1.4305 & 1.4305 & 1.000 & \checkmark \\
  501192 & bin & 18 & 8 & 0 & 1.4048 & 1.4390 & 1.024 & \checkmark \\
  501281 & bin & 35 & 9 & 3 & 1.3563 & 1.3591 & 1.002 & \checkmark \\
  481195 & cont & 7 & 5 & 0 & 1.2116 & 1.2116 & 1.000 & \checkmark \\
  501429 & cont & 74 & 7 & 9 & 1.1104 & 1.1104 & 1.000 & \checkmark \\
  501330 & bin & 29 & 7 & 0 & 1.0047 & 0.00e+00 & --- & $\times$ \\
  482630 & cont & 26 & 17 & 2 & 0.8973 & 0.8973 & 1.000 & \checkmark \\
  480804 & bin & 19 & 6 & 0 & 0.8466 & 0.8312 & 0.982 & \checkmark \\
  479993 & cont & 73 & 35 & 12 & 0.8376 & 0.8376 & 1.000 & \checkmark \\
  480926 & bin & 9 & 8 & 1 & 0.8153 & 0.8153 & 1.000 & \checkmark \\
  473552 & bin & 32 & 9 & 0 & 0.7897 & 0.7873 & 0.997 & \checkmark \\
  501392 & bin & 51 & 12 & 2 & 0.6928 & 0.6937 & 1.001 & \checkmark \\
  501355 & bin & 11 & 9 & 1 & 0.6668 & 0.6668 & 1.000 & \checkmark \\
  501226 & bin & 12 & 5 & 0 & 0.6181 & 0.00e+00 & --- & $\times$ \\
  501209 & bin & 20 & 11 & 4 & 0.5991 & 0.5991 & 1.000 & \checkmark \\
  501337 & bin & 20 & 4 & 1 & 0.5964 & 0.5964 & 1.000 & \checkmark \\
  501384 & bin & 60 & 15 & 3 & 0.5782 & 0.00e+00 & --- & $\times$ \\
  479767 & bin & 12 & 5 & 0 & 0.5645 & --- & --- & --- \\
  501324 & bin & 13 & 4 & 3 & 0.5482 & 0.5482 & 1.000 & \checkmark \\
  501194 & bin & 94 & 17 & 8 & 0.5362 & 0.5188 & 0.968 & \checkmark \\
  501409 & cont & 51 & 8 & 7 & 0.5236 & 0.5236 & 1.000 & \checkmark \\
  481384 & bin & 16 & 4 & 0 & 0.4544 & 0.00e+00 & --- & $\times$ \\
  501253 & bin & 13 & 11 & 5 & 0.4344 & 0.4323 & 0.995 & \checkmark \\
  482362 & bin & 7 & 5 & 2 & 0.4303 & 0.4303 & 1.000 & \checkmark \\
  501268 & bin & 37 & 5 & 0 & 0.4227 & 0.3898 & 0.922 & \checkmark \\
  479600 & bin & 30 & 9 & 1 & 0.4119 & 0.4119 & 1.000 & \checkmark \\
  480492 & bin & 40 & 7 & 3 & 0.4117 & 0.4062 & 0.987 & \checkmark \\
  480935 & bin & 7 & 5 & 1 & 0.4087 & 0.3404 & 0.833 & $\times$ \\
  482001 & bin & 102 & 7 & 3 & 0.3690 & 0.00e+00 & --- & $\times$ \\
  480629 & bin & 23 & 10 & 0 & 0.3607 & 0.3607 & 1.000 & \checkmark \\
  501235 & bin & 20 & 7 & 0 & 0.3550 & 0.00e+00 & --- & $\times$ \\
  481856 & bin & 22 & 4 & 1 & 0.3497 & 0.00e+00 & --- & $\times$ \\
  481140 & cont & 21 & 4 & 1 & 0.3488 & 0.3843 & 1.102 & $\times$ \\
  481193 & bin & 45 & 4 & 1 & 0.3423 & 0.3423 & 1.000 & \checkmark \\
  479806 & bin & 7 & 5 & 0 & 0.3085 & 0.3085 & 1.000 & \checkmark \\
  479969 & bin & 95 & 21 & 1 & 0.3082 & 0.00e+00 & --- & $\times$ \\
  480430 & bin & 15 & 4 & 0 & 0.3041 & 0.2820 & 0.927 & \checkmark \\
  501412 & bin & 14 & 6 & 0 & 0.3033 & 0.3033 & 1.000 & \checkmark \\
  501305 & bin & 12 & 5 & 0 & 0.2875 & 0.2875 & 1.000 & \checkmark \\
  501395 & bin & 14 & 8 & 4 & 0.2736 & 0.2736 & 1.000 & \checkmark \\
  501367 & bin & 19 & 6 & 2 & 0.2724 & 0.2634 & 0.967 & \checkmark \\
  481412 & cont & 25 & 18 & 11 & 0.2698 & 0.2699 & 1.000 & \checkmark \\
  479808 & bin & 27 & 7 & 0 & 0.2643 & 0.00e+00 & --- & $\times$ \\
  501393 & bin & 27 & 7 & 0 & 0.2643 & 0.2644 & 1.000 & \checkmark \\
  501287 & bin & 13 & 7 & 5 & 0.2591 & 0.2231 & 0.861 & $\times$ \\
  501388 & bin & 34 & 14 & 5 & 0.2589 & 0.2566 & 0.991 & \checkmark \\
  479974 & bin & 13 & 4 & 0 & 0.2427 & 0.00e+00 & --- & $\times$ \\
  481378 & bin & 93 & 13 & 6 & 0.2192 & --- & --- & --- \\
  481974 & bin & 11 & 6 & 0 & 0.2157 & 0.2157 & 1.000 & \checkmark \\
  481589 & bin & 41 & 10 & 0 & 0.2139 & 0.2139 & 1.000 & \checkmark \\
  501416 & cont & 24 & 9 & 5 & 0.2125 & 0.2125 & 1.000 & \checkmark \\
  501277 & cont & 13 & 4 & 1 & 0.2076 & 0.2076 & 1.000 & \checkmark \\
  501261 & bin & 21 & 10 & 0 & 0.2056 & 0.2057 & 1.000 & \checkmark \\
  482476 & bin & 15 & 9 & 1 & 0.2053 & 0.00e+00 & --- & $\times$ \\
  481988 & bin & 7 & 6 & 2 & 0.1918 & 0.0223 & 0.116 & $\times$ \\
  481107 & bin & 30 & 5 & 0 & 0.1864 & 0.1857 & 0.996 & \checkmark \\
  501424 & bin & 43 & 9 & 2 & 0.1841 & 0.1841 & 1.000 & \checkmark \\
  481009 & bin & 32 & 13 & 0 & 0.1823 & 0.1824 & 1.000 & \checkmark \\
  501348 & bin & 89 & 4 & 3 & 0.1767 & 0.1766 & 0.999 & \checkmark \\
  501283 & bin & 55 & 17 & 2 & 0.1760 & 0.1760 & 1.000 & \checkmark \\
  482440 & bin & 16 & 5 & 2 & 0.1755 & 0.1755 & 1.000 & \checkmark \\
  481736 & cont & 37 & 4 & 0 & 0.1692 & 0.1692 & 1.000 & \checkmark \\
  479585 & bin & 48 & 10 & 0 & 0.1654 & 0.00e+00 & --- & $\times$ \\
  482734 & bin & 48 & 4 & 1 & 0.1622 & 0.1538 & 0.948 & \checkmark \\
  481588 & bin & 26 & 18 & 3 & 0.1606 & 0.1606 & 1.000 & \checkmark \\
  501214 & bin & 41 & 11 & 1 & 0.1489 & 0.1488 & 0.999 & \checkmark \\
  482450 & bin & 11 & 4 & 0 & 0.1420 & 0.1420 & 1.000 & \checkmark \\
  482519 & bin & 11 & 7 & 4 & 0.1346 & 0.1346 & 1.000 & \checkmark \\
  481323 & bin & 13 & 9 & 3 & 0.1330 & 0.1330 & 1.000 & \checkmark \\
  501399 & cont & 29 & 4 & 1 & 0.1299 & 0.1299 & 1.000 & \checkmark \\
  480851 & cont & 29 & 4 & 1 & 0.1290 & 0.1291 & 1.000 & \checkmark \\
  482004 & bin & 25 & 8 & 2 & 0.1251 & 0.1243 & 0.994 & \checkmark \\
  481216 & bin & 18 & 5 & 1 & 0.1248 & 0.1248 & 1.000 & \checkmark \\
  501215 & bin & 10 & 7 & 2 & 0.1216 & 0.1216 & 1.000 & \checkmark \\
  501387 & bin & 15 & 5 & 1 & 0.1183 & 0.1184 & 1.000 & \checkmark \\
  480060 & bin & 46 & 5 & 1 & 0.1126 & 0.1126 & 1.000 & \checkmark \\
  482382 & cont & 9 & 5 & 0 & 0.1116 & 0.1118 & 1.002 & \checkmark \\
  481733 & bin & 12 & 4 & 2 & 0.1038 & 0.1038 & 1.000 & \checkmark \\
  480706 & bin & 8 & 5 & 3 & 0.1023 & 0.1023 & 1.000 & \checkmark \\
  501390 & bin & 26 & 8 & 2 & 0.0972 & 0.0969 & 0.997 & \checkmark \\
  501263 & bin & 9 & 5 & 0 & 0.0959 & 0.0960 & 1.001 & \checkmark \\
  481470 & cont & 12 & 6 & 0 & 0.0944 & 0.0944 & 1.000 & \checkmark \\
  482472 & bin & 5 & 4 & 0 & 0.0926 & 0.0926 & 1.001 & \checkmark \\
  501377 & bin & 17 & 8 & 5 & 0.0805 & 0.0804 & 0.999 & \checkmark \\
  501313 & bin & 19 & 8 & 0 & 0.0711 & 0.0711 & 1.000 & \checkmark \\
  501435 & bin & 69 & 7 & 0 & 0.0684 & 0.0684 & 1.000 & \checkmark \\
  481418 & bin & 12 & 4 & 0 & 0.0678 & 0.0679 & 1.001 & \checkmark \\
  501318 & bin & 17 & 7 & 1 & 0.0672 & 0.0673 & 1.002 & \checkmark \\
  501382 & bin & 17 & 11 & 2 & 0.0666 & 0.0665 & 0.998 & \checkmark \\
  481781 & bin & 6 & 5 & 0 & 0.0649 & 0.0642 & 0.988 & \checkmark \\
  480655 & cont & 12 & 9 & 5 & 0.0624 & 0.0624 & 1.000 & \checkmark \\
  501272 & bin & 64 & 14 & 12 & 0.0610 & 0.0610 & 1.000 & \checkmark \\
  481695 & bin & 11 & 6 & 0 & 0.0568 & 0.0568 & 1.001 & \checkmark \\
  479907 & bin & 59 & 11 & 1 & 0.0558 & 0.0558 & 1.000 & \checkmark \\
  479770 & bin & 23 & 4 & 0 & 0.0554 & 0.0172 & 0.311 & $\times$ \\
  481908 & bin & 16 & 4 & 0 & 0.0502 & 0.0503 & 1.001 & \checkmark \\
  480667 & bin & 59 & 9 & 7 & 0.0452 & 0.0452 & 1.000 & \checkmark \\
  501338 & bin & 7 & 5 & 1 & 0.0419 & 0.0419 & 1.000 & \checkmark \\
  482003 & bin & 8 & 5 & 2 & 0.0360 & 0.0360 & 1.001 & \checkmark \\
  480695 & bin & 145 & 5 & 0 & 0.0356 & 0.00e+00 & --- & $\times$ \\
  482159 & bin & 30 & 6 & 8 & 0.0339 & 0.0339 & 1.001 & \checkmark \\
  501311 & bin & 14 & 5 & 0 & 0.0318 & 0.0319 & 1.003 & \checkmark \\
  501256 & bin & 61 & 5 & 3 & 0.0311 & 0.0311 & 1.000 & \checkmark \\
  473269 & bin & 14 & 5 & 0 & 0.0281 & 0.0282 & 1.005 & \checkmark \\
  480666 & bin & 45 & 29 & 8 & 0.0266 & 0.0259 & 0.973 & \checkmark \\
  481147 & bin & 5 & 4 & 0 & 0.0261 & 0.0263 & 1.007 & \checkmark \\
  501346 & bin & 58 & 7 & 0 & 0.0250 & 0.0249 & 0.995 & \checkmark \\
  480612 & cont & 23 & 5 & 2 & 0.0235 & 0.0235 & 0.999 & \checkmark \\
  481731 & bin & 14 & 6 & 0 & 0.0233 & 0.0234 & 1.003 & \checkmark \\
  481200 & bin & 5 & 4 & 2 & 0.0229 & 0.0229 & 1.002 & \checkmark \\
  501250 & bin & 19 & 8 & 4 & 0.0226 & 0.0226 & 1.002 & \checkmark \\
  479531 & bin & 22 & 19 & 2 & 0.0215 & 0.0216 & 1.002 & \checkmark \\
  482258 & bin & 29 & 4 & 1 & 0.0209 & 0.0209 & 1.000 & \checkmark \\
  501423 & cont & 10 & 4 & 1 & 0.0203 & 0.0203 & 1.001 & \checkmark \\
  501366 & cont & 20 & 7 & 2 & 0.0202 & 0.0202 & 1.000 & \checkmark \\
  481118 & bin & 13 & 8 & 0 & 0.0191 & 0.0191 & 1.002 & \checkmark \\
  501228 & bin & 46 & 9 & 1 & 0.0187 & 0.0187 & 0.999 & \checkmark \\
  501321 & cont & 39 & 9 & 2 & 0.0139 & 0.0139 & 1.003 & \checkmark \\
  481734 & bin & 28 & 6 & 3 & 0.0138 & 0.0139 & 1.004 & \checkmark \\
  501257 & bin & 22 & 6 & 4 & 0.0136 & 0.0137 & 1.004 & \checkmark \\
  501232 & bin & 111 & 12 & 2 & 0.0132 & 0.0133 & 1.008 & \checkmark \\
  481150 & bin & 31 & 11 & 0 & 0.0096 & 0.0107 & 1.110 & $\times$ \\
  481163 & bin & 10 & 8 & 0 & 0.0085 & 0.0143 & 1.677 & $\times$ \\
  501376 & bin & 87 & 17 & 0 & 0.0079 & 0.0081 & 1.022 & \checkmark \\
  501373 & bin & 28 & 7 & 4 & 0.0068 & 0.0069 & 1.011 & \checkmark \\
  482109 & bin & 14 & 6 & 0 & 0.0067 & 0.0068 & 1.014 & \checkmark \\
  481236 & bin & 6 & 4 & 0 & 0.0059 & 0.0070 & 1.187 & $\times$ \\
  479773 & bin & 11 & 5 & 0 & 0.0053 & 0.0055 & 1.029 & \checkmark \\
  482465 & bin & 15 & 10 & 4 & 0.0047 & 0.0061 & 1.290 & $\times$ \\
  501374 & bin & 34 & 5 & 7 & 0.0045 & 0.0046 & 1.028 & \checkmark \\
  501420 & bin & 10 & 5 & 0 & 0.0045 & 0.0050 & 1.122 & $\times$ \\
  501282 & cont & 16 & 9 & 1 & 0.0042 & 0.0042 & 0.997 & \checkmark \\
  481263 & cont & 11 & 7 & 1 & 0.0040 & 0.0040 & 0.989 & \checkmark \\
  501345 & bin & 18 & 5 & 0 & 0.0033 & 0.0040 & 1.207 & $\times$ \\
  501206 & bin & 49 & 8 & 4 & 0.0030 & 0.0031 & 1.031 & \checkmark \\
  501224 & bin & 11 & 9 & 1 & 0.0018 & 0.0071 & 3.971 & $\times$ \\
  481836 & cont & 16 & 9 & 11 & 4.00e-04 & 4.47e-04 & 1.116 & \checkmark \\
  501408 & cont & 27 & 6 & 0 & 2.00e-04 & 3.54e-04 & 1.772 & $\times$ \\
  501430 & bin & 19 & 10 & 3 & 3.17e-09 & 0.0036 & --- & $\times$ \\
  501298 & bin & 6 & 5 & 0 & 3.05e-09 & --- & --- & --- \\
  501370 & bin & 10 & 5 & 0 & 2.24e-09 & 0.0176 & --- & $\times$ \\
  501372 & cont & 9 & 5 & 1 & 1.78e-09 & 0.0086 & --- & $\times$ \\
  501326 & bin & 25 & 4 & 0 & 1.08e-09 & 0.0028 & --- & $\times$ \\
  480563 & bin & 8 & 5 & 0 & 9.89e-10 & 0.0102 & --- & $\times$ \\
  481551 & bin & 11 & 6 & 1 & 8.77e-10 & 0.0086 & --- & $\times$ \\
  481942 & bin & 41 & 6 & 7 & 7.25e-10 & 0.00e+00 & --- & \checkmark \\
  474842 & bin & 45 & 8 & 2 & 4.62e-10 & 0.0036 & --- & $\times$ \\
  482576 & bin & 11 & 5 & 0 & 4.54e-10 & 0.0049 & --- & $\times$ \\
  501407 & bin & 62 & 4 & 0 & 4.14e-10 & 0.0041 & --- & $\times$ \\
  501297 & bin & 52 & 4 & 6 & 3.87e-10 & 0.0040 & --- & $\times$ \\
  480074 & bin & 16 & 5 & 0 & 3.77e-10 & 0.0051 & --- & $\times$ \\
  501434 & bin & 14 & 6 & 3 & 3.70e-10 & 0.0039 & --- & $\times$ \\
  480931 & bin & 13 & 6 & 1 & 3.67e-10 & 0.0035 & --- & $\times$ \\
  481766 & bin & 8 & 7 & 6 & 3.02e-10 & 0.0041 & --- & $\times$ \\
  479996 & bin & 21 & 9 & 2 & 2.96e-10 & 0.0047 & --- & $\times$ \\
  501227 & bin & 21 & 8 & 0 & 1.94e-10 & 0.0015 & --- & $\times$ \\
  501336 & bin & 8 & 4 & 0 & 1.85e-10 & 0.0035 & --- & $\times$ \\
  501312 & bin & 22 & 9 & 3 & 1.83e-10 & --- & --- & --- \\
  501293 & bin & 19 & 4 & 1 & 1.68e-10 & 0.0023 & --- & $\times$ \\
  480001 & bin & 10 & 7 & 3 & 1.59e-10 & 0.0036 & --- & $\times$ \\
  480052 & bin & 20 & 8 & 2 & 1.49e-10 & 0.0024 & --- & $\times$ \\
  482416 & bin & 8 & 7 & 0 & 1.41e-10 & --- & --- & --- \\
  501340 & bin & 20 & 4 & 0 & 1.04e-10 & 0.0018 & --- & $\times$ \\
  501207 & bin & 22 & 8 & 1 & 9.88e-11 & 0.0021 & --- & $\times$ \\
  479661 & bin & 14 & 4 & 2 & 9.40e-11 & 0.0016 & --- & $\times$ \\
  501300 & cont & 12 & 4 & 1 & 8.80e-11 & 0.0025 & --- & $\times$ \\
  479807 & bin & 14 & 5 & 0 & 8.28e-11 & 0.0019 & --- & $\times$ \\
  501365 & bin & 13 & 4 & 0 & 8.27e-11 & 0.0014 & --- & $\times$ \\
  501243 & bin & 17 & 7 & 0 & 7.97e-11 & 0.0016 & --- & $\times$ \\
  501197 & bin & 11 & 5 & 2 & 7.83e-11 & 0.0015 & --- & $\times$ \\
  501255 & bin & 29 & 4 & 0 & 7.76e-11 & --- & --- & --- \\
  501350 & bin & 40 & 11 & 3 & 7.68e-11 & 0.0013 & --- & $\times$ \\
  501309 & bin & 11 & 5 & 2 & 6.63e-11 & 0.0013 & --- & $\times$ \\
  482120 & bin & 16 & 9 & 0 & 5.78e-11 & 0.0018 & --- & $\times$ \\
  479574 & bin & 18 & 7 & 0 & 5.16e-11 & 0.0010 & --- & $\times$ \\
  501325 & bin & 26 & 5 & 0 & 4.86e-11 & 9.43e-04 & --- & $\times$ \\
  481571 & bin & 10 & 8 & 1 & 4.45e-11 & 0.0014 & --- & $\times$ \\
  482451 & cont & 11 & 8 & 0 & 3.85e-11 & 0.0023 & --- & $\times$ \\
  501343 & bin & 14 & 5 & 0 & 2.98e-11 & 8.68e-04 & --- & $\times$ \\
  479971 & bin & 27 & 4 & 0 & 2.35e-11 & 5.37e-04 & --- & $\times$ \\
  479650 & bin & 27 & 6 & 1 & 1.87e-11 & 6.79e-04 & --- & $\times$ \\
  480037 & bin & 11 & 7 & 0 & 1.53e-11 & 0.0010 & --- & $\times$ \\
  501252 & bin & 14 & 7 & 0 & 1.25e-11 & 8.51e-04 & --- & $\times$ \\
  501419 & bin & 7 & 6 & 0 & 1.22e-11 & 0.0022 & --- & $\times$ \\
  482522 & bin & 6 & 4 & 0 & 9.01e-12 & 8.46e-04 & --- & $\times$ \\
  482477 & cont & 5 & 4 & 0 & 6.80e-12 & 4.08e-04 & --- & $\times$ \\
  476033 & bin & 7 & 5 & 0 & 6.14e-12 & 4.84e-04 & --- & $\times$ \\
  501302 & bin & 8 & 6 & 0 & 6.14e-12 & 5.17e-04 & --- & $\times$ \\
  481159 & bin & 8 & 5 & 0 & 5.38e-12 & 4.93e-04 & --- & $\times$ \\
  501381 & bin & 12 & 10 & 6 & 4.52e-12 & 7.48e-04 & --- & $\times$ \\
  501267 & bin & 26 & 9 & 5 & 3.98e-12 & 2.36e-04 & --- & $\times$ \\
  482521 & bin & 6 & 4 & 0 & 3.05e-12 & 5.68e-04 & --- & $\times$ \\
  501404 & bin & 23 & 5 & 3 & 2.88e-12 & 3.19e-04 & --- & $\times$ \\
  480029 & bin & 17 & 8 & 0 & 2.24e-12 & 7.68e-04 & --- & $\times$ \\
  501204 & bin & 27 & 7 & 0 & 1.41e-12 & 5.20e-04 & --- & $\times$ \\
  482231 & bin & 7 & 4 & 0 & 3.50e-13 & 6.90e-04 & --- & $\times$ \\
  479918 & bin & 15 & 7 & 0 & 2.49e-13 & 8.39e-04 & --- & $\times$ \\
  482527 & bin & 10 & 5 & 0 & 2.33e-13 & 0.0011 & --- & $\times$ \\
\bottomrule
\end{longtable}

% Summary: 128 OK, 84 fail, 0 err out of 218

\endgroup

\end{document}
